\section{Theoretical Framework and Approach}
\subsection{Approach}

This article develops theory and does not test it. Its design corresponds to what~\citet{jaakkola2020designing} calls a model paper, in which the author specifies constructs and the relationships among them and argues for why those relationships should hold. Following~\citet{whetten1989constitutes}, we try to be explicit about which factors are included, how they connect, and what logic links them. Because the central claims concern how an outcome varies across conditions, we state them as propositions, the format that~\citet{cornelissen2017developing} associates with theorizing about contingent relationships. Throughout the paper, a claim attributed to a cited study reports what that study found or argued. Claims that are ours are flagged as such, and none of them has been tested.

\subsection{Social Learning Theory}

Social learning theory holds that people acquire much of what they do by watching others and noting what happens to them, as well as by acting and experiencing consequences directly~\citep{bandura1977social}. Bandura later broadened the account into a social cognitive theory in which personal factors, behavior, and the environment influence one another~\citep{bandura1986social}. Applied to organizations, observational learning has been described as running through attention to a model, retention of what was seen, reproduction of the behavior, and motivation to act on it~\citep{manz1981vicarious}. Early field work found that subordinates' behavior resembled that of their superiors more closely when the superiors were seen as competent and successful~\citep{weiss1977subordinate}, and the same author extended the argument from behavior to work values~\citep{weiss1978social}.

The theory is also the main foundation of research on ethical leadership. \citet{brown2005ethical} defined ethical leaders partly as credible role models whose conduct followers observe and emulate, and whose communication, rewards, and discipline signal which conduct is expected. A later review drew heavily on social learning reasoning to explain why ethical leadership affects followers~\citep{brown2006ethical}. Bandura also recognized that people learn from verbal and symbolic models, such as descriptions and depictions of conduct, and not only from live ones~\citep{bandura1977social}. Recent theory adds that vicarious learning need not require silent observation: it can also happen in conversation, when people jointly work through what one of them experienced~\citep{myers2018coactive}.

In our reading, each of these steps can be affected by distance. Attention requires that a model be present to attend to, and retention and reproduction depend on having seen a behavior clearly enough to remember and repeat it. Motivation depends partly on what learners see happen to the model afterwards. Bandura called this last element vicarious reinforcement: observing that others are rewarded or punished for a behavior informs a learner's expectations about what will follow if the learner acts the same way~\citep{bandura1977social}. We suggest that for ethical norms vicarious reinforcement is often the decisive element, because what an organization rewards and punishes in practice is what distinguishes its enacted norms from its stated ones.

Two features of the theory make it fit our problem. It locates the learning of norms in exposure to models and to the consequences that follow their conduct, so any change in exposure should change what is learned. And it treats the consequences others experience as information, which is precisely the kind of information that tells a newcomer which ethical norms are enforced rather than merely stated.

\subsection{Imprinting Theory}

The idea of imprinting entered organization theory through~\citet{stinchcombe1965social}, who observed that organizations bear lasting marks of the period in which they were founded. \citet{marquis2013imprinting} generalized the idea across levels of analysis: during a brief sensitive period, an entity develops characteristics that reflect prominent features of its environment, and those characteristics persist after the environment changes. They also separated imprinting from related ideas, including path dependence and cohort effects.

At the level of individuals, research has linked early career conditions to lasting patterns. A longitudinal study of professionals found that the economic conditions of their organization at the time of hire shaped the approaches they developed during socialization and, through those approaches, their later performance~\citep{tilcsik2014imprint}. Young scientists have been shown to adopt their advisers' orientation toward commercial science~\citep{azoulay2017social}. Organizations have been described as leaving career imprints that employees carry into later roles~\citep{higgins2005career}, and prior work experience has been found to carry over in ways that can hinder performance in a new setting~\citep{dokko2009unpacking}.

The concept of a sensitive period is central to this literature. Imprinting is not expected to happen at any moment, but during periods in which an entity is unusually open to its environment, and periods of transition, such as founding or entry into a new setting, are commonly discussed examples~\citep{marquis2013imprinting}. In our view, entry into a first professional job has several such features for an individual: the newcomer has few established routines, is actively trying to learn what is expected, and has little prior experience against which to compare what the organization presents.

Imprinting theory fits our problem because it explains what social learning theory leaves open: why learning that happens at entry should matter after entry. It directs attention to the conditions present during the sensitive period, and it predicts that later exposure will not fully overwrite what was formed then.

\subsection{Why These Two Theories}

The two theories answer different questions, which is why we use both. Social learning theory explains how enacted ethical norms are acquired and why the channel of exposure matters. Imprinting theory explains why the result of that acquisition, including its gaps, persists. We considered and set aside several alternatives. Sensemaking accounts of entry~\citep{louis1980surprise} explain how newcomers interpret surprises but not why early interpretations endure. Research on ethical climate and culture describes the content that newcomers learn, and we draw on it for that purpose, but it does not specify the learning process. Identity-based accounts of generations~\citep{joshi2010unpacking} help us define which cohort matters, but they are not a theory of how ethical norms are transmitted. A control or monitoring perspective would predict that distance increases opportunities for misconduct; we treat it as a rival explanation and return to it in the discussion.

\subsection{Scope Conditions}

The model applies to office-based and knowledge work, where distributed arrangements are feasible, and to the entry period in a new organization, with entry into a first professional job as the clearest case. It concerns the ethical norms of the work unit a newcomer joins, especially in the gray areas where written rules leave room for judgment and enacted norms therefore carry most weight; that emphasis is our assumption. It does not address roles that cannot be performed remotely, and it takes no position on whether distributed work is good or bad overall.
